\chapter{Metodología}
La metodología elegida tiene un impacto fundamental en el desarrollo del proyecto,
existen diferentes y variadas opciones, personalmente no fue hasta que comencé a 
indagar en este tema que me comencé a sentir seguro de mi trabajo, realizar 
cualquier proyecto era más cuestión de suerte y ganas que de puesta en 
practica del conocimiento. De las muchas opciones, me decanto por dos:
\begin{description}
	\item[RUP] Rational Unified Process.
	\item[XP] Extreme Programming.
\end{description}
\section{Dos es mejor que una}
¿Se puede desarrollar un proyecto con dos metodologías?, la respuesta depende del
tipo de proyecto, en este caso la respuesta es un contundente si, el proyecto es
muy amplio implica partes mecánicas, eléctricas y de programación, mi propuesta 
es utilizar \emph{RUP} para las partes mecánicas y eléctricas y \emph{XP} para 
la programación.
\subsection{RUP}
Dentro de las metodologías hay dos polos, las ágiles y las guiadas por plan, 
RUP
es una metodología guiada por plan, sin embargo no está en un extremo opuesto a
las ágiles.

RUP consiste en una serie de pasos bien definidos, esta diseñada para abarcar 
proyectos muy grandes, proyectos que implican colaboraciones entre países o entre
empresas transnacionales, mi proyecto es pequeño en comparación, sin embargo 
considero adecuado aprovechar varias de las características de esta metodología,
los siguientes escalones son el esqueleto del proceso de desarrollo:
\begin{description}
	\item[Requisitación] En este primer paso estamos en el peor de los mundos, pues
		tenemos poca información del problema y una especificación pobre.
		Hacer una toma de requisitos en la metodología RUP, implica 
		entender el lenguaje del dominio, no es lo mismo hablar con un
		médico que con
		un agricultor o un chef de cocina, entender a quien tiene el 
		problema en es un paso fundamental para poder diseñar
		una solución que satisfaga sus necesidades, no entender el 
		lenguaje implicaría una serie de adivinanzas desafortunadas.

		Cuando el entendimiento del lenguaje es suficiente, la prioridad
		es encontrar los requerimientos, nuevamente los requerimientos 
		vienen de parte de quien requiere la solución.
	\item[Análisis] Con los requerimientos definidos, procedemos a analizarlos,
		el proceso de análisis consiste en transformar los
		requerimientos del lenguaje del dominio a el lenguaje de la 
		ingeniería, cada requerimiento 
		puede abarcar varias áreas, para el robot móvil implica 
		mecánica, electronica y programación. 
	\item[Diseño] Cada área de la ingeniería tiene una forma de diseño propia,
		en el caso de la mecánica y la electronica consiste en crear la 
		arquitectura que satisfaga el requerimiento y posteriormente 
		hacer los cálculos de los parámetros que la
		conforman. En el caso de la programación implica definir la 
		arquitectura y encontrar las entidades que la satisfagan, la 
		granularidad propia del software dificulta encontrar tales 
		entidades. 
	\item[Implemetación] La implementación es el trabajo pesado y manual, 
		aquí es donde se ensamblan los componentes, se alambran los 
		circuitos y se escribe el código, es la etapa más lenta de todas,
		el punto de hacer una metodología es establecer las condiciones
		necesarias
		para que esta etapa pueda fluir lo más suave posible.
	\item[Pruebas] Realizada la implementación se debe verificar que cumpla 
		con las expectativas, si no lo hicieran debe repetirse el proceso
		nuevamente, no se trata de iniciar desde cero, si no de corregir
		lo que se erró en el proceso.
\end{description}

La intensión es realizar esta cascada de pasos para cada uno de los requisitos,
las veces que sean necesarias hasta que se alcancen los mismos.
En las primeras iteraciones la requisición se llevaran la mayor parte del tiempo,
y con forme avancen las iteraciones irán perdiendo protagonismo pues cada vez 
estarán más refinados, en las ultimas iteraciones las pruebas se llevarán el
protagonismo.

\subsection{XP}
La metodología \emph{Extreme Programming} (XP), es una de las revolucionarias
de los últimos
tiempos, a diferencia de de RUP, donde se separan las fases de desarrollo, en 
XP el análisis, diseño, pruebas e implementación se van haciendo al mismo tiempo,
XP es el padre de las modernas SCRUM y KANBAN entre muchas otras, está diseñada
para proyectos medianos y pequeños, pero se requiere una habilidad técnica muy alta
para llévala acabo además de que solo se recomienda para proyectos de software.

La técnica insignia de las metodologías ágiles es el diseño guiado por pruebas
TDD\cite{tdd}, Test Drive Development en ingles, esta técnica consiste en 
escribir la prueba primero y después 
escribir solo el código necesario para que la prueba pase, la premisa del TDD
es: \emph{la mejor forma de corregir bugs es evitar introducirlos}.
Al realizar las pruebas primero se obliga a escribir código correcto, gracias a 
las multiples herramientas existentes sobre pruebas unitarias, es posible 
automatizar el proceso, personalmente he adoptado el TDD y ya no
concibo programar de otra forma, aunque reconozco que proceso fue traumático,
tuve que literalmente olvidar lo aprendido y volver a aprender con nuevas
directivas.

\section{¿Porqué no usar una metodología clásica?}
Hay varias razones. El proceso de desarrollo de un proyecto según el PmBok es
el siguiente
\begin{enumerate}
	\item Inicio.- Definir el proyecto.
	\item Planificación.- Se establecen los objetivos, se analizan y se
		diseña la solución.
	\item Implementación
	\item Medición de desempeño.
	\item Cierre.
\end{enumerate}

Ciertamente no difiere mucho de RUP, sin embargo no implica iteración en ningún
momento, definir el proyecto al inicio es muy difícil cuando de proyectos 
innovadores se trata, lo más probable es que vamos a fallar, pienso que es más
valioso aprender a lidiar con los cambios que planificar un proceso perfecto. 
El diseño temprano e implementación tardía implica una rigidez enorme, si
se añadieran requisitos obligarían a rediseñar todo el proyecto. ¿Y la planificación?, 
¿acaso no es bueno planificar?, ami parecer planificar a corto plazo es bueno, 
planificar grandes cantidades de trabajo es un ejercicio inútil.

--Vamos a diseñar este proyecto y en dos meses el 5 de Junio comenzamos a 
implementarlo.--

Permítame reírme de sus planes.

¿Qué pasa si una vez implementado el proyecto no pasa las pruebas?, con suerte 
los cambios para pasar las pruebas serán lo suficientemente pequeños para poder
desplegarlos, pero si no es así, si por alguna razón se nos escapo un detalle
en el diseño o si el cliente añade una nueva característica, entonces sera el
momento adecuado para hechar a andar nuestra imaginación para inventar las
escusas más creativas, pues implementar los cambios requerirán un proceso 
lento y doloroso.

El mundo ideal no existe, la experiencia me indica que las cosas siempre 
cambian, una metodología flexible es la que decido utilizar.  
